\documentclass[11pt]{article}

\usepackage[a4paper,margin=0.82in]{geometry}
\usepackage[T1]{fontenc}
\usepackage[utf8]{inputenc}
\usepackage{lmodern}
\usepackage{amsmath,amssymb,amsthm,mathtools,mathrsfs}
\usepackage{booktabs,enumitem,microtype,tabularx,xcolor}
\usepackage[numbers,sort&compress]{natbib}
\usepackage[colorlinks=true,linkcolor=blue!65!black,
  citecolor=blue!65!black,urlcolor=blue!70!black]{hyperref}
\usepackage[nameinlink,noabbrev]{cleveref}
\hypersetup{
  pdftitle={A Provenance-Preserving Quotient of Literal Affine Maps},
  pdfauthor={Liang Wang},
  pdfsubject={TPC-105 exact map quotient and collision splitting},
  pdfkeywords={affine map, provenance, collision energy, fixed shift, resonance carrier}
}

\newtheorem{theorem}{Theorem}[section]
\newtheorem{proposition}[theorem]{Proposition}
\newtheorem{corollary}[theorem]{Corollary}
\newtheorem{lemma}[theorem]{Lemma}
\theoremstyle{definition}
\newtheorem{definition}[theorem]{Definition}
\newtheorem{example}[theorem]{Example}
\theoremstyle{remark}
\newtheorem{remark}[theorem]{Remark}

\newcommand{\F}{\mathbb F}
\newcommand{\Z}{\mathbb Z}
\newcommand{\one}{\mathbf 1}
\newcommand{\Lzero}{\mathrm{L0}}
\newcommand{\Lone}{\mathrm{L1}}
\newcommand{\Ltwo}{\mathrm{L2}}
\newcommand{\PX}{\mathfrak P_X}
\newcommand{\XX}{\mathfrak X_X}
\newcommand{\KX}{\mathfrak K_X}
\newcommand{\YX}{\mathfrak Y_X}

\title{\textbf{A Provenance-Preserving Quotient}\\
\textbf{of Literal Affine Maps:}\\
Exact Collision Splitting and the Identical-Map Concentration Ledger}
\author{Liang Wang\\
School of Artificial Intelligence and Automation\\
Huazhong University of Science and Technology, Wuhan 430074, P.R. China\\
\texttt{wangliang.f@gmail.com}}
\date{July 2026}

\begin{document}
\maketitle

\begin{abstract}
TPC-100 linearizes every remaining invertible opened atom as
\[
 \omega_p=A_p(m_pj_p+h_0)
 \quad\text{in }\mathbb F_{q_X}^{\times},
\]
but its maximally literal fine cells are singletons.  Thus all
distinct-atom collisions remain in the cross-cell ledger, and
different cells may induce the same affine map.  We construct the
canonical quotient required before a genuine cross-map estimate is
attempted.

Attach to an atom \(p\) the map
\[
 \phi_p(x)=s_px+t_p,\qquad
 s_p=A_pm_p,\quad t_p=A_ph_0.
\]
Atoms are equivalent when their pairs \((s_p,t_p)\) agree.  A
lossless quotient record consists of the full recoverable
provenance bag together with the input profile
\[
 \nu_\phi(j)=
 \sum_{\substack{p:\phi_p=\phi\\j_p=j}}w_p.
\]
The profile, not merely the total map weight, is necessary and
sufficient for the positive multiplier census:
\[
 a_H^{\mathrm{inv}}(\omega)
 =
 \sum_\phi\sum_j
 \nu_\phi(j)\one_{\{\phi(j)=\omega\}}.
\]
We give a two-point counterexample proving that map totals alone
lose the census.

Because \(h_0,A_p,m_p\) are nonzero modulo the strong no-wrap prime,
two maps agree exactly when \(A_p=A_{p'}\) and
\(m_p=m_{p'}\) modulo \(q_X\); no-wrap upgrades the latter to
equality of the physical row integers.  Within one map, output
collision is exactly equality of the input coordinates.
Consequently the TPC-100 centered cross ledger splits identically as
\[
 \mathfrak X_H^\circ
 =
 \mathfrak K_H^\circ+\mathfrak Y_H^\circ,
\]
where \(\mathfrak K_H^\circ\) contains identical maps and
\(\mathfrak Y_H^\circ\) contains genuinely distinct maps.

If \(M_{H,\phi}\) is the mass of one map class,
\(U_X=\max_{H,\phi}M_{H,\phi}\), and \(M_H\) is the total width
mass, then
\[
 |\mathfrak K_H^\circ|\le U_XM_H,
\]
and therefore
\[
 \sum_Hg_{q_X}(H)\sqrt{|\mathfrak K_H^\circ|}
 \le
 \sqrt{U_X|\mathcal H_X|(q_X-1)\mathfrak P_X}.
\]
This is an exact \(\Lone\) reduction.  It exposes a new
identical-map concentration ledger \(U_X\) and leaves the
width-weighted distinct-map term as the honest TPC-106 geometric
door.  Neither ledger is bounded here at the growing physical
scale; no \(\Ltwo\) arithmetic saving or parity breakthrough is
claimed.
\end{abstract}

\noindent\textbf{Keywords:}
affine map; provenance quotient; collision energy; identical maps;
incidence; fixed shift; positive resonance.

\medskip
\noindent\textbf{Literal-frame convention.}
Every quotient starts from the actual \(q_X\nmid u\) opened atoms,
keeps one prescribed \(h_0\ne0\), their full source and child
records, actual weights and support, both polarizations, and the
single inherited physical normalization.  Nothing is specialized
to \(h_0=2\).

\section{Why a quotient is needed}
\label{sec:introduction}

Fix an exact resonance width \(H\).  TPC-100 partitions the actual
invertible opened atoms into maximally literal fine cells.  Because
each cell retains \(m,u,\sigma\), the source identity
\[
 mj+h_0=\sigma u
\]
recovers \(j\); every fine cell is therefore a singleton.  The
resulting multiplier census is
\begin{equation}
 a_H^{\mathrm{inv}}(\omega)
 =
 \sum_{\substack{p:\,\mathsf H_p=H\\q_X\nmid u_p}}
 w_p\one_{\{\omega_p=\omega\}},
\label{eq:census}
\end{equation}
where \(w_p\ge0\) is the exact physical atom weight
\citep{WangTPC100}.

The centered energy has an atomic diagonal part and a cross part
\[
 \|(a_H^{\mathrm{inv}})^\circ\|_2^2
 =
 \mathfrak D_H^\circ+\mathfrak X_H^\circ.
\]
TPC-103 closes the maximum atom cap, so the diagonal part follows
from the unresolved principal mass
\citep{WangTPC101,WangTPC103}.  The cross part remains.  Before
applying an incidence theorem to it, TPC-102 and TPC-MVP1 require
identical affine maps to be aggregated without losing source
recovery or exact weight \citep{WangTPC102,WangTPCMVP1}.

The subtle point is that an affine map is an operator on an input,
not merely a multiplier value.  Two atoms may carry the same map
but occur at different \(j\)'s.  Their total class mass does not
determine where that mass is sent.

\section{Literal affine maps}
\label{sec:maps}

Write \(q=q_X\).  For all sufficiently large \(X\), the strong
no-wrap choice gives
\[
 q>|h_0|,\qquad 0<m_p<q,\qquad q\nmid u_p.
\]
TPC-100 defines
\[
 A_p=
 \varepsilon_p\ell_pv_pB_p(c_pu_p)^{-1}
 \in\F_q^\times
\]
and proves
\begin{equation}
 \omega_p=A_p(m_pj_p+h_0).
\label{eq:linearization}
\end{equation}

\begin{definition}[Literal affine-map record]
\label{def:map}
For an active atom \(p\), define
\[
 s_p=A_pm_p,\qquad t_p=A_ph_0,\qquad
 \phi_p(x)=s_px+t_p.
\]
Thus \(s_p,t_p\in\F_q^\times\) and
\(\omega_p=\phi_p(j_p)\).
\end{definition}

\begin{lemma}[Exact identical-map classification]
\label{lem:classification}
For two active atoms \(p,p'\),
\[
 \phi_p=\phi_{p'}
 \quad\Longleftrightarrow\quad
 A_p=A_{p'}\ \text{and}\ m_p=m_{p'}
 \quad\text{in }\F_q.
\]
Under the strong no-wrap bounds, equality of \(m_p,m_{p'}\) in
\(\F_q\) is equality as physical integers.
\end{lemma}

\begin{proof}
Equality of the constant terms gives
\[
 h_0(A_p-A_{p'})=0.
\]
Since \(q\nmid h_0\), this forces \(A_p=A_{p'}\).  Equality of the
slopes and invertibility of \(A_p\) then force
\(m_p=m_{p'}\).  The converse is immediate.  Finally both physical
row integers lie strictly between \(0\) and \(q\)
\citep{WangTPC98}.
\end{proof}

\begin{corollary}[Collision inside one map]
\label{cor:same-map-collision}
If \(\phi_p=\phi_{p'}\), then
\[
 \omega_p=\omega_{p'}
 \quad\Longleftrightarrow\quad
 j_p=j_{p'}\pmod q.
\]
Since the actual orbit window has diameter below \(q\), this is
equivalent to equality of the physical integers \(j_p=j_{p'}\).
\end{corollary}

\begin{proof}
The common slope \(s_p\) is nonzero, so
\(\phi_p(j_p)=\phi_p(j_{p'})\) exactly when
\(j_p=j_{p'}\) modulo \(q\).  Apply no-wrap.
\end{proof}

\begin{remark}[What identical maps may still hide]
Equality of \((A,m,j)\) does not identify the full atom.  Different
content or projector children, complementary factorizations, masks,
polarizations, or outer source labels may still contribute separate
weights.  Those records may not be discarded.
\end{remark}

\section{The canonical provenance quotient}
\label{sec:quotient}

Let \(\mathscr P_H\) be the finite set of active invertible atoms at
width \(H\).  Define
\[
 p\sim p'
 \quad\Longleftrightarrow\quad
 (s_p,t_p)=(s_{p'},t_{p'}).
\]
Let \(\Phi_H\) be the set of equivalence classes, identified with
their common affine maps.

\begin{definition}[Provenance bag and input profile]
\label{def:profile}
For \(\phi\in\Phi_H\), let
\[
 \mathscr B_\phi
 =
 \{(p,w_p,\mathfrak r_p):\phi_p=\phi\},
\]
where \(\mathfrak r_p\) is the complete TPC-100 augmented source
and child record.  Define
\begin{align}
 \nu_\phi(j)
 &=
 \sum_{\substack{p\in\mathscr P_H\\
                  \phi_p=\phi,\ j_p=j}}w_p,
\label{eq:profile}\\
 M_{H,\phi}
 &=
 \sum_j\nu_\phi(j),\qquad
 M_H=\sum_{\phi\in\Phi_H}M_{H,\phi}.
\label{eq:class-mass}
\end{align}
The canonical quotient record is
\[
 \mathscr Q_H
 =
 \{(\phi,\nu_\phi,\mathscr B_\phi):\phi\in\Phi_H\}.
\]
\end{definition}

\begin{theorem}[Lossless quotient and exact census]
\label{thm:quotient}
The classes \(\mathscr B_\phi\) form a disjoint partition of the
actual atoms, preserve every physical weight exactly once, and give
\begin{equation}
 \boxed{
 a_H^{\mathrm{inv}}(\omega)
 =
 \sum_{\phi\in\Phi_H}\sum_{j\in\F_q}
 \nu_\phi(j)\one_{\{\phi(j)=\omega\}}.}
\label{eq:quotient-census}
\end{equation}
The analytic data \((\phi,\nu_\phi)\) recover the complete positive
multiplier census, while the attached bags recover every original
source atom and coefficient.
\end{theorem}

\begin{proof}
Every atom has one and only one pair \((s_p,t_p)\), so the bags form
a partition.  Reindex \eqref{eq:census} first by that pair and then
by \(j_p\).  This gives \eqref{eq:quotient-census} and
\[
 \sum_{\phi,j}\nu_\phi(j)=\sum_pw_p=M_H.
\]
No weight is created or deleted.  Source recovery follows because
the full augmented record \(\mathfrak r_p\) remains in the
corresponding bag.
\end{proof}

\begin{proposition}[Map totals alone are not lossless]
\label{prop:mass-only}
The collection \(\{(\phi,M_{H,\phi})\}\) does not determine
\(a_H^{\mathrm{inv}}\), even for one invertible affine map and one
unit atom.
\end{proposition}

\begin{proof}
Take \(q\ge5\), \(h_0=1\), and
\(\phi(x)=x+1\).  One unit atom at \(j=0\) gives the census
\(\delta_1\), while one unit atom at \(j=1\) gives
\(\delta_2\).  Both quotient records have the same map and total
mass one.  Their input profiles, and hence their multiplier
censuses, differ.
\end{proof}

\begin{remark}[Minimal analytic payload]
The proposition shows why ``aggregate identical maps with exact
total weight'' is incomplete unless the input distribution is also
retained.  The profile \(\nu_\phi\) is sufficient for every positive
incidence calculation.  The provenance bag is additionally required
for literal source recovery and later signed reconstruction.
\end{remark}

\section{Exact splitting of the cross ledger}
\label{sec:splitting}

Since the TPC-100 fine cells are singleton, its centered cross
ledger can be written directly as
\begin{equation}
 \mathfrak X_H^\circ
 =
 \sum_{\substack{p,p'\in\mathscr P_H\\p\ne p'}}
 w_pw_{p'}
 \left(
 \one_{\{\omega_p=\omega_{p'}\}}-\frac1{q-1}
 \right).
\label{eq:cross}
\end{equation}
The sum is ordered, matching the \(C\ne D\) convention in TPC-100.

\begin{definition}[Identical- and distinct-map ledgers]
\label{def:split}
Let
\begin{align}
 \mathfrak K_H^\circ
 &=
 \sum_{\substack{p\ne p'\\\phi_p=\phi_{p'}}}
 w_pw_{p'}
 \left(
 \one_{\{\omega_p=\omega_{p'}\}}-\frac1{q-1}
 \right),
\label{eq:K}\\
 \mathfrak Y_H^\circ
 &=
 \sum_{\substack{p,p'\\\phi_p\ne\phi_{p'}}}
 w_pw_{p'}
 \left(
 \one_{\{\omega_p=\omega_{p'}\}}-\frac1{q-1}
 \right).
\label{eq:Y}
\end{align}
\end{definition}

\begin{theorem}[Canonical collision split]
\label{thm:split}
One has the exact identity
\begin{equation}
 \boxed{
 \mathfrak X_H^\circ
 =
 \mathfrak K_H^\circ+\mathfrak Y_H^\circ.}
\label{eq:split}
\end{equation}
Moreover, put
\[
 S_{\phi,j}
 =
 \sum_{\substack{p:\phi_p=\phi\\j_p=j}}w_p^2,
\qquad
 S_\phi=\sum_jS_{\phi,j}.
\]
Then the identical-map term has the profile formula
\begin{equation}
 \boxed{
 \mathfrak K_H^\circ
 =
 \sum_{\phi,j}
 \left(\nu_\phi(j)^2-S_{\phi,j}\right)
 -
 \frac1{q-1}
 \sum_\phi
 \left(M_{H,\phi}^2-S_\phi\right).}
\label{eq:K-profile}
\end{equation}
The distinct-map term is
\begin{equation}
 \boxed{
 \mathfrak Y_H^\circ
 =
 \sum_{\substack{\phi,\psi\in\Phi_H\\\phi\ne\psi}}
 \sum_{j,k}
 \nu_\phi(j)\nu_\psi(k)
 \left(
 \one_{\{\phi(j)=\psi(k)\}}-\frac1{q-1}
 \right).}
\label{eq:Y-profile}
\end{equation}
\end{theorem}

\begin{proof}
The alternatives \(\phi_p=\phi_{p'}\) and
\(\phi_p\ne\phi_{p'}\) partition every ordered pair \(p\ne p'\),
giving \eqref{eq:split}.  Within one map,
\cref{cor:same-map-collision} replaces output equality by
\(j_p=j_{p'}\).  At fixed \((\phi,j)\), the mass of ordered distinct
pairs is
\[
 \nu_\phi(j)^2-S_{\phi,j}.
\]
Across all inputs of the same map, the total ordered distinct-pair
mass is
\[
 M_{H,\phi}^2-S_\phi.
\]
This proves \eqref{eq:K-profile}.  Reindexing pairs from two
different classes by their two input coordinates gives
\eqref{eq:Y-profile}.
\end{proof}

\begin{remark}[Centering may couple the two terms]
Neither \(\mathfrak K_H^\circ\) nor
\(\mathfrak Y_H^\circ\) is necessarily nonnegative.  Their exact
sum is the native cross excess.  Bounding their absolute values
separately is a sufficient route and may lose cancellation; it is
not asserted to be necessary.
\end{remark}

\section{The identical-map concentration ledger}
\label{sec:concentration}

\begin{definition}[Maximum map-class mass]
\label{def:U}
Set
\[
 U_H=\max_{\phi\in\Phi_H}M_{H,\phi},
\qquad
 U_X=\max_HU_H,
\]
with zero used for an empty carrier.
\end{definition}

\begin{theorem}[Identical-map concentration bound]
\label{thm:K-bound}
For every width \(H\),
\begin{equation}
 \boxed{
 |\mathfrak K_H^\circ|
 \le
 \sum_{\phi\in\Phi_H}
 \left(M_{H,\phi}^2-S_\phi\right)
 \le U_HM_H.}
\label{eq:K-bound}
\end{equation}
\end{theorem}

\begin{proof}
For one ordered pair, the centered collision kernel is either
\[
 1-\frac1{q-1}
 \quad\text{or}\quad
 -\frac1{q-1},
\]
so its modulus is at most one.  Summing the weights of all ordered
distinct pairs inside identical-map classes gives the first bound.
For the second,
\[
 \sum_\phi(M_{H,\phi}^2-S_\phi)
 \le\sum_\phi M_{H,\phi}^2
 \le U_H\sum_\phi M_{H,\phi}
 =U_HM_H.
\]
\end{proof}

Let
\[
 \mathcal H_X=
 \{H:g_{q_X}(H)>0,\ M_H>0\},
\qquad
 \PX=\sum_H\frac{2H}{q_X-1}M_H.
\]

\begin{corollary}[Width-weighted transfer]
\label{cor:transfer}
The identical-map ledger obeys
\begin{equation}
 \boxed{
 \sum_Hg_{q_X}(H)
 \sqrt{|\mathfrak K_H^\circ|}
 \le
 \sqrt{
 U_X|\mathcal H_X|(q_X-1)\PX}.}
\label{eq:transfer}
\end{equation}
Hence, if
\[
 q_X\asymp Q_X,\qquad
 U_X\le X^{o(1)},\qquad
 \PX\le X^{o(1)}Q_X^2,
\]
then the left side of \eqref{eq:transfer} is
\(X^{o(1)}Q_X^2\).
\end{corollary}

\begin{proof}
By \cref{thm:K-bound},
\[
 \sum_Hg_q(H)\sqrt{|\mathfrak K_H^\circ|}
 \le
 \sqrt{U_X}\sum_{H\in\mathcal H_X}
 g_q(H)\sqrt{M_H}.
\]
Cauchy--Schwarz and \(g_q(H)^2\le2H\) give
\[
 \le
 \sqrt{
 U_X|\mathcal H_X|
 \sum_H2HM_H}
 =
 \sqrt{
 U_X|\mathcal H_X|(q-1)\PX}.
\]
Use \(|\mathcal H_X|\le(q-1)/2\).
\end{proof}

\begin{corollary}[Exact sufficient distinct-map reduction]
\label{cor:distinct}
Suppose the hypotheses in \cref{cor:transfer} hold and
\begin{equation}
 \sum_Hg_{q_X}(H)
 \sqrt{|\mathfrak Y_H^\circ|}
 \le X^{o(1)}Q_X^2.
\label{eq:distinct-target}
\end{equation}
Then the complete width-weighted cross gate satisfies
\[
 \sum_Hg_{q_X}(H)
 \sqrt{|\mathfrak X_H^\circ|}
 \le X^{o(1)}Q_X^2.
\]
\end{corollary}

\begin{proof}
Use \(\mathfrak X_H^\circ=\mathfrak K_H^\circ+
\mathfrak Y_H^\circ\) and
\[
 \sqrt{|x+y|}\le\sqrt{|x|}+\sqrt{|y|}.
\]
Then apply \cref{cor:transfer} and \eqref{eq:distinct-target}.
\end{proof}

\section{Two exact nonimplications}
\label{sec:nonimplications}

\begin{proposition}[Atom cap does not imply map-class cap]
\label{prop:atom-vs-map}
For every \(R\ge1\), there is an abstract positive opened census
with
\[
 W_X=1,\qquad U_X=R.
\]
It can be supported on one invertible affine map.
\end{proposition}

\begin{proof}
Take \(R\) distinct provenance records, each of weight one, and
attach all of them to the same map \(\phi(x)=x+1\).  Their map-class
mass is \(R\), while every atom has weight one.
\end{proof}

\begin{proposition}[A quotient does not create dispersion]
\label{prop:no-dispersion}
Canonical aggregation preserves the census and its centered energy
exactly.  In particular, no bound for
\(\mathfrak Y_H^\circ\) or \(U_X\) follows from the quotient
identity alone.
\end{proposition}

\begin{proof}
Equation \eqref{eq:quotient-census} is an equality with the original
census.  Therefore its \(\ell^2\) norm and centered energy are
unchanged.  \Cref{prop:atom-vs-map} shows that one quotient class
may have arbitrary positive mass in the abstract information
model.
\end{proof}

\begin{remark}[Physical status]
The countermodels prove nonimplications between abstract ledgers.
They do not establish polynomial concentration in the actual TPC
coefficient and hence do not stop the physical positive route.  A
future theorem may use the exact classification in
\cref{lem:classification} to bound actual class multiplicity or
weight.
\end{remark}

\section{Finite exact regression}
\label{sec:certificate}

The deterministic script
\texttt{experiments/tpc105\_map\_quotient\_audit.py} constructs
finite weighted atom families over several prime fields and checks:
\begin{enumerate}[leftmargin=2em]
\item direct and quotient multiplier censuses agree;
\item the map classification in \cref{lem:classification};
\item the direct cross ledger equals both
\(\mathfrak K_H^\circ+\mathfrak Y_H^\circ\) and the profile
formulas;
\item the concentration inequality \eqref{eq:K-bound}; and
\item the mass-only and atom-cap counterexamples.
\end{enumerate}
All arithmetic is exact.  The certificate tests the finite
reindexing and claim schema; it is not evidence for a growing
distinct-map incidence estimate.

\section{Route consequence and claim boundary}
\label{sec:boundary}

\subsection*{Established}

\begin{itemize}[leftmargin=2em]
\item Identical affine maps are classified exactly by the pair
\((A_p,m_p)\), with physical integer recovery for \(m_p\).
\item The canonical quotient retains the input profile needed for
the positive census and the full bag needed for source recovery.
\item The TPC-100 cross ledger splits exactly into identical- and
distinct-map terms.
\item The identical-map term is controlled by the new class-mass
ledger \(U_X\), and \cref{cor:distinct} gives a precise sufficient
door for TPC-106.
\end{itemize}

These are \(\Lzero\) algebraic identities and an \(\Lone\)
provenance-preserving reduction on the actual fixed-\(h_0\) carrier.

\subsection*{Not established}

\begin{itemize}[leftmargin=2em]
\item No bound \(U_X\le X^{o(1)}\) is proved for the actual carrier.
\item The principal bounds isolated in TPC-104 remain unproved
\citep{WangTPC104}.
\item The distinct-map target \eqref{eq:distinct-target} is not
proved.
\item No signed affine M\"obius cancellation, full physical
reassembly, strict \(1/400\) endpoint closure, or \(\Ltwo\)
fixed-power saving is obtained.
\item No parity breakthrough, Hardy--Littlewood asymptotic,
prime-pair lower bound, or twin-prime conclusion is proved or
implied.
\end{itemize}

The honest next step is width-weighted incidence on the genuine
distinct-map quotient, with \(U_X\) carried as a separate
concentration gate.  If either has a literal polynomial obstruction,
the response is to retain the signed filter before majorization,
not to claim that the original signed coefficient is large.

\begingroup
\scriptsize
\setlength{\bibsep}{1pt}
\raggedright
\bibliographystyle{plainnat}
\bibliography{references}
\endgroup

\end{document}
